% Test-case specification TC-0012 (CTest test TC-0012-regcall).
\documentclass[letterpaper]{article}

\usepackage[letterpaper,margin=1in]{geometry}
\emergencystretch=3em

% Compile from the repository root so this path resolves to
% the pinned WSP submodule.
\input{wsp/testing/test-case-library.tex}

\begin{document}

\def\TCID{TC-0012}
\def\TCTitle{%
  Helpers called with operands in registers; Zombie interface}
\def\TCPurpose{%
  Verify the helpers called as compiled code calls them, with their operands
  in registers (the Integer, Byte, Abs and Sgn helpers, the string, object
  and Variant free, copy and move helpers), and the Zombie interface.}
\def\TCPriority{\TCPriorityHigh}

\def\TCRequirementRef{%
  REQ-0002, REQ-0003}
\def\TCDesignRef{%
  Not applicable.}
\def\TCFeatureRef{%
  \texttt{tests/\allowbreak{}vbrun\_\allowbreak{}regcall\_\allowbreak{}tests.\allowbreak{}cpp},
  \texttt{src/\allowbreak{}helpers\_\allowbreak{}regcall.\allowbreak{}cpp}}
\def\TCTestTechnique{%
  Requirements-based testing of the calling convention.}

\def\TCPreconditions{%
  The library and the test program are built
  (\texttt{scripts/\allowbreak{}Invoke-Build.\allowbreak{}ps1}) and the test
  is registered with CTest.}
\def\TCEnvironment{%
  Windows with the toolchain of
  \texttt{scripts/\allowbreak{}Invoke-Build.\allowbreak{}ps1}; the local
  gate builds and runs x64 Debug, CI builds and runs the other
  architectures. The library is built from this repository in
  \texttt{out/\allowbreak{}build}.}
\def\TCAssumptions{%
  Not applicable.}

\def\TCInputData{%
  Operands loaded into registers by the test's call stubs.}
\def\TCInitialState{%
  No state from earlier runs is needed; the test makes what it needs.}

\def\TCProcedure{%
  \begin{enumerate}
    \item CTest runs \texttt{vbrun\_\allowbreak{}regcall\_\allowbreak{}tests}.
    \item Call each helper through a stub that loads the registers the compiler
      uses.
    \item Check the results, the registers and the stack after the call.
    \item CTest records the exit code and, on failure, the program's output
      (\texttt{--output-on-failure}).
  \end{enumerate}}
\def\TCExpectedResult{%
  All checks pass; the program prints \texttt{vbrun regcall: all checks
  passed} and exits 0.}
\def\TCPostConditions{%
  No file or registry state remains that later tests depend on.}

\TCTemplate

\end{document}
